\chapter{Predict pictures from reconstructed samples}
\MARSChapterHeading{Predict pictures from reconstructed samples}
\label{revision:predicting-pictures}

The lossless video format described changes with XOR. Its decoder recovered every
original sample exactly. H.263 takes another approach: predict a region from an
earlier picture, transform the remaining difference, then choose how much of that
information to retain. The reconstructed picture can differ from the source.
That loss is useful only if the decoder still interprets the same stream and
both sides retain compatible prediction state.

We will build a restricted baseline H.263 encoder across this chapter and the
next. Here we follow samples through comparison, prediction and local
reconstruction. The next chapter describes the selected coding decisions as
bits and writes a stream for an independent decoder. The complete worked source
is \texttt{h263-encoder.asm} in \path{revision/labs/files}. The larger
\texttt{035-h263} package retains the callable contracts and construction work.
Use those contracts when implementing a stage; use this chapter to understand
why the stages connect as they do.

The input is raw QCIF video generated by Python, limited to thirty pictures per
second and ten seconds. Our ordinary example is much smaller: eight pictures at
$5000/1001$ pictures per second, lasting 1.6016 seconds including the last hold.
It contains a stationary background and a moving square. We already know MSA;
the challenge is choosing correct lane widths and memory windows while preserving
an interoperable decoder's meaning.

\section{Locate a sample in three separate planes}

A raw picture has no header announcing its dimensions or pixel format. The
accompanying description supplies that information. This input contains a
$176\times144$ Y plane followed by $88\times72$ Cb and Cr planes. Y records
luma; Cb and Cr record color differences at half the horizontal and vertical
sample density. This arrangement is called planar 4:2:0. It is not an RGB word
array or an interleaved sequence of three channel bytes.

\begin{table}[!ht]
  \centering
  \begin{tabular}{l r r r r}
    \toprule
    Plane & Byte offset & Width & Height & Row stride \\
    \midrule
    Y     & 0           & 176   & 144    & 176        \\
    Cb    & 25344       & 88    & 72     & 88         \\
    Cr    & 31680       & 88    & 72     & 88         \\
    \bottomrule
  \end{tabular}
  \caption{The tightly packed planes occupy 38016 bytes per raw picture.}
\end{table}

Picture $f$ starts at file offset $38016f$. Its Y sample at $(x,y)$ is a further
$176y+x$ bytes into that picture. Thus picture one's sample at $(16,32)$ is at
$38016+32\cdot176+16=43664$. The corresponding chroma origin is $(8,16)$;
use stride 88 and the appropriate plane offset there. A byte sample has the same
numerical meaning in both target byte orders. Signed halfword residuals and
native request fields introduced later do depend on byte order.

The generator keeps its nominal Y values within 16--235 and chroma within
16--240. The numerical kernels still accept all sample bytes from zero through
255. Do not replace a callable range with the narrower values that happen to
occur in one example. Likewise, the 12:11 display sample aspect ratio does not
change the stored width or stride. A display may stretch the picture horizontally;
the encoder still reads the same 176 luma samples per row.

A sixteen-by-sixteen luma macroblock contains four eight-by-eight blocks. Its
corresponding Cb and Cr regions each supply one more block, for six blocks in
all. Number the Y blocks top-left, top-right, bottom-left, bottom-right, then
Cb and Cr. The shared block size lets prediction, transforms and coefficient
coding use small interfaces without confusing one plane's coordinates with
another's.

\section{Measure a candidate before selecting motion}

For current samples $C$ and a candidate reference window $R$, define the sum of
absolute differences, or SAD:
\[
  S=\sum_{y=0}^{15}\sum_{x=0}^{15}|C[y,x]-R[y,x]|.
\]
Each magnitude is at most 255, so the complete score is at most
$256\cdot255=65280$. A sixteen-by-sixteen window of 130 compared with 128 has
score $256\cdot2=512$. No signed byte interpretation belongs in this subtraction:
byte 255 means a bright sample, not minus one.

The scalar and MSA SAD entries take two top-left pointers and two independent
strides. Each row contributes sixteen bytes, even when the stride is 176 or
4096. A pointer can be unaligned. Equal pointers do not prove equal windows if
the strides differ. Both inputs remain read-only, including when they overlap.
Admit strides and non-wrapping address extents before reading samples.

In the vector path, one sixteen-byte load fits the active row exactly. Absolute
unsigned subtraction gives sixteen byte magnitudes. Widen adjacent pairs before
accumulating across rows:
\Needspace{9\baselineskip}
\MARSCodeLabel{Accumulate one row of absolute sample differences}
\begin{minted}{mars-idea}
    ld.b $w0, 0($t0)
    ld.b $w1, 0($t1)
    asub_u.b $w2, $w0, $w1
    hadd_u.h $w3, $w2, $w2
    addv.h $w4, $w4, $w3
\end{minted}
Here \texttt{w4} starts at zero before the row loop. Each halfword receives two
magnitudes per row, at most 510, and at most $16\cdot510=8160$ over the block.
Later horizontal additions widen again before forming the full score. Keeping
the row accumulation in bytes would discard carries even when the final result
is stored in a wide general register.

Motion search keeps the current window fixed and compares reference windows at
$(x+d_x,y+d_y)$. The worked policy searches full-pixel displacements in a radius
zero, one or two. An interior radius-two search has 25 candidates. At the upper
left corner only offsets zero through two survive in each dimension, leaving
nine. Clip the candidate intervals before constructing addresses; do not clamp
several out-of-picture candidates onto the same edge sample.

For a complete reference block, its origin must satisfy
$0\leq x+d_x\leq160$ and $0\leq y+d_y\leq128$. Apply the reference stride to the
admitted row, then add its admitted column. The current pointer already addresses
the current block; adding $(x,y)$ to it again would move the window twice.
Positive displacement selects a position to the right or below in the reference.
An object moving down and right can therefore need a negative reference vector.

A score alone does not resolve ties. The project prefers zero motion if it has
the minimum score, otherwise increasing $d_y$ and then increasing $d_x$.
A centre-first traversal that replaces only for a strictly smaller score gives
this order. A reverse traversal can compare the same tie key explicitly.
These are encoder choices. H.263 does not require this search radius or tie rule,
and SAD does not account for the bits needed to code a vector or its residual.

\begin{figure}[!ht]
\centering
\begin{tikzpicture}[alt={The native caller observation record contains success zero and SAD 7872 for two generated luma windows. The score field is outlined.}]
\node[anchor=north west,inner sep=0,font=\small\bfseries] (title) at (0,0) {Read a complete block score};
\node[anchor=north west,inner sep=0] (a) at ([yshift=-2mm]title.south west) {\MARSRevisionCrop{revision/captures/h263-prediction/h263-sad.png}{0.99\linewidth}{0.03125}{0.33888889}{0.58203125}{0.60888889}};
\begin{scope}[shift={(a.south west)},x={($(a.south east)-(a.south west)$)},y={($(a.north west)-(a.south west)$)}]
\draw[black,line width=.7pt] (0.582500,0.005000) rectangle (0.965000,0.443809);
\end{scope}
\end{tikzpicture}
\caption{The caller copies success and SAD into two native doublewords. These generated raw windows give 7872 (outlined); they isolate the comparison kernel before reconstructed-reference prediction.}
\label{fig:revision-h263-native-sad}
\end{figure}


\section{Interpolate the reference that the decoder will use}

The search returns whole luma pixels, but H.263 represents motion at half-pixel
precision. The later syntax stage must convert units. Chroma also has lower
sample density: a displacement of one luma pixel corresponds to half a chroma
sample. Whole-pixel luma search therefore still requires fractional chroma
prediction.

For a signed luma component $d$, split its chroma displacement into
\[
  q=\lfloor d/2\rfloor,\qquad p=d-2q.
\]
The integer origin shift is $q$ and the phase $p$ is zero or one. With $d=-1$,
we need origin minus one and phase one, giving minus one half. Division that
truncates toward zero would choose the wrong origin. An arithmetic right shift
forms the floor, and the low bit gives the phase. At chroma position $(8,16)$,
luma vector $(-1,-1)$ consequently selects origin $(7,15)$ and both half phases.

Let $A$ be the integer reference sample, $B$ its right neighbour, $C$ the one
below, and $D$ the below-right neighbour. The selected baseline rounding gives
\[
  \begin{array}{ll}
    P_{0,0}=A,                       & P_{1,0}=\lfloor(A+B+1)/2\rfloor, \\[2pt]
    P_{0,1}=\lfloor(A+C+1)/2\rfloor, &
    P_{1,1}=\lfloor(A+B+C+D+2)/4\rfloor.
  \end{array}
\]
Sum all four diagonal neighbours before rounding. If $A=B=C=0$ and $D=1$,
the diagonal prediction is zero. Rounding each row average upward and then
rounding their average upward would instead produce one. The same interpolation
formula must be used in the encoder's local decoder and the receiver.

\begin{figure}[!ht]
\centering
\begin{tikzpicture}[every node/.style={font=\small},alt={A square of reference samples has top-left, top-right and bottom-left zero, bottom-right one. The centre half-sample is floor of one plus two divided by four, equal to zero.}]
\draw[black!50] (0,0) rectangle (4,2);
\draw[dashed,black!50] (0,1)--(4,1) (2,0)--(2,2);
\node[fill=white] at (0,2) {$A=0$};\node[fill=white] at (4,2) {$B=0$};
\node[fill=white] at (0,0) {$C=0$};\node[fill=white] at (4,0) {$D=1$};
\node[fill=white,draw,align=center] at (2,1) {half sample\\$P=0$};
\end{tikzpicture}
\caption{The diagonal prediction rounds the complete four-neighbour sum once.}
\label{fig:revision-h263-interpolation}
\end{figure}


An eight-by-eight prediction needs an $(8+p_x)$-by-$(8+p_y)$ reference footprint.
Diagonal prediction uses nine rows and nine columns, while the current block
still uses eight of each. The extra row and column are neighbours, not arbitrary
padding. An MSA load of sixteen bytes from an eight-byte current row is therefore
not automatically valid. Gather the active samples into halfword lanes, or prove
that a larger load contains only admitted samples.

The four-neighbour sum plus bias is at most 1022, so unsigned halfwords suffice.
After predicting, widen current samples and form signed residuals $E=C-P$ in
$[-255,255]$. The prediction entry returns both arrays: sixty-four predicted
bytes followed by sixty-four native signed halfwords. A residual SAD summarizes
magnitudes, but does not replace those signed corrections. For the generated
Cb example, current 128 and neighbours 128,128,128,96 give prediction 120 and
residual eight. The complete block's residual SAD is 232.

\begin{figure}[!ht]
  \centering
  \begin{tikzpicture}[alt={The observation record contains success zero, block residual SAD 232, first predicted sample 120 and first signed residual eight. Prediction is outlined and residual has a dashed outline.}]
    \node[anchor=north west,inner sep=0,font=\small\bfseries] (title) at (0,0) {Keep prediction and signed residual distinct};
    \node[anchor=north west,inner sep=0] (a) at ([yshift=-2mm]title.south west) {\MARSRevisionCrop{revision/captures/h263-prediction/h263-prediction.png}{0.99\linewidth}{0.03125}{0.31555556}{0.58203125}{0.60888889}};
    \begin{scope}[shift={(a.south west)},x={($(a.south east)-(a.south west)$)},y={($(a.north west)-(a.south west)$)}]
      \draw[black,line width=.7pt] (0.200000,0.005000) rectangle (0.582500,0.305824);
      \draw[black,dashed,line width=.7pt] (0.582500,0.005000) rectangle (0.965000,0.305824);
    \end{scope}
  \end{tikzpicture}
  \caption{The caller copies status, block SAD, first predicted sample and first signed residual. The outlined 120 and dashed eight come from the generated chroma window with both half phases; 232 summarizes the entire residual block.}
  \label{fig:revision-h263-native-prediction}
\end{figure}


\section{Express the residual as frequency coefficients}

An eight-by-eight discrete cosine transform, or DCT, changes the block's basis.
Instead of storing a value at each spatial position, it stores weights for
sixty-four horizontal/vertical patterns. The DC coefficient describes the constant
component; AC coefficients describe variation. There are still sixty-four numbers.
The transform alone has not compressed them.

Use row-first notation $f[y,x]$ for samples and $F[v,u]$ for coefficients, where
$v$ is vertical frequency and $u$ horizontal frequency. Define
\[
  B[k,n]=\gamma_k\cos\!\left(\frac{(2n+1)k\pi}{16}\right),\quad
  \gamma_0=1/\sqrt8,\quad \gamma_k=1/2\ (k>0).
\]
The forward transform is
\[
  F[v,u]=\sum_{y=0}^{7}\sum_{x=0}^{7}B[v,y]f[y,x]B[u,x].
\]
One pass combines each row with horizontal patterns; the other combines those
results vertically. A constant block of value two has DC sixteen and zero AC.
The factor eight follows from the normalization in both dimensions. A spatial
impulse or an asymmetric ramp helps detect an exchanged row or column that a
constant block would conceal.

For INTER coding, transform the signed prediction residual. For INTRA coding,
transform absolute sample values directly; do not subtract 128 because another
image codec used that convention. A constant INTRA block of 128 has DC 1024.
The inverse transform uses the transposed basis in each dimension. Keeping the
request's direction separate from its mode prevents scalar/MSA selection from
accidentally selecting a different mathematical operation.

The worked implementation stores a basis scaled by $2^{16}$. First-pass sums
retain that fractional scale in signed words. The second pass carries scale
$2^{32}$ in signed doublewords and rounds once at the end. Rounding the first
pass back to samples discards information before the vertical combination.
Transposing the private intermediate array makes each second-pass input contiguous;
that storage decision does not change the transform's orientation.

Scale, product width and rounding bias must agree. A different fixed-point
precision can be accurate, but substituting a smaller basis while retaining the
old shift will change the answer. The callable forward boundary clips coefficients
to $[-2048,2047]$; the inverse clips spatial results to $[-256,255]$. The latter
are signed corrections for INTER work, not yet display bytes. Finite examples
allow the specified one-unit transform accuracy. The separate Annex A test also
checks statistical error limits and exact zero reconstruction across its prescribed
cohorts; a round trip through the implementation's own forward transform would
not supply an independent inverse check.

\section{Distinguish selected levels from reconstructed coefficients}

Quantization maps several possible forward coefficients to one selected integer
level. The decoder receives the level, so it cannot recover which original
coefficient produced it. For the worked forward policy and $1\leq Q\leq31$, use
\[
  m=\min(127,\lfloor |C|/(2Q)\rfloor),\qquad L=\operatorname{sign}(C)m.
\]
At $Q=8$, magnitudes below sixteen disappear. Coefficients sixteen and minus
sixteen become levels one and minus one. This zero region reduces what later
needs coding. The cap also matters: at very small $Q$, a large coefficient can
hit level 127. Smaller $Q$ alone does not guarantee arbitrarily close reconstruction.

Forward selection is our reproducible policy. The decoder's inverse meaning is
fixed. Zero ordinary level reconstructs zero. For nonzero $L$, form magnitude
$Q(2|L|+1)$, subtract one if $Q$ is even, restore the sign, then clip to
$[-2048,2047]$. At $Q=8$, levels one and minus one reconstruct as 23 and minus
23. Applying the nonzero formula to level zero would manufacture a correction;
subtracting the even-parameter adjustment after restoring the sign would treat
negative coefficients differently.

Baseline INTRA DC has a separate representation. Select an integer $d$ in
one through 254 at eight-unit coefficient spacing; the worked policy rounds
$(C+4)/8$ downward and clamps to that range. Code $d=128$ as byte 255.
The decoder interprets that special code as coefficient 1024. Ordinary $Q$
quantization still applies to the other 63 coefficients. An INTRA DC code is
therefore neither the original coefficient nor an ordinary signed AC level.

Keep selected levels and inverse-quantized coefficients in distinct arrays.
Coefficient coding will consume the former; the local inverse transform consumes
the latter. Both can use native halfword storage without having interchangeable
meanings. A valid complete encoder can choose other legal levels, provided its
stream, quality and local reconstruction satisfy the full assignment. The numerical
exercise's exact policy is not a requirement for identical whole-file bytes.

\section{Close the reconstruction loop before predicting again}

Suppose an original picture sample was 100 but reconstructed as 98. The next raw
sample is 104. Subtracting the old original gives residual four; adding it to the
decoder's retained 98 gives 102. The subtraction and addition were individually
correct. They used different reference pictures. Prediction must read the previous
accepted reconstruction on both sides.

\begin{figure}[!ht]
  \centering
  \begin{tikzpicture}[node distance=5mm,every node/.style={draw,align=center,font=\small,text width=93mm},alt={Current minus prediction forms signed residuals. Forward transform and quantization select levels. Those levels supply stream coding and the local decoder. Inverse quantization and inverse transform produce signed corrections. Adding the same prediction and clipping produces the next reference picture.}]
    \node (a) {Current minus prediction: signed residuals};
    \node[below=of a] (b) {Forward transform and quantization: selected levels};
    \node[below=of b] (c) {Inverse quantization and transform: signed correction};
    \node[below=of c] (d) {Add the same prediction and clip: next reference};
    \draw[-{Latex}] (a)--(b);\draw[-{Latex}] (b)--(c);\draw[-{Latex}] (c)--(d);
  \end{tikzpicture}
  \caption{The local decoder starts from the selected levels that are also coded into the stream. It does not undo quantization using the original coefficients.}
  \label{fig:revision-h263-reconstruction-loop}
\end{figure}


In the actual block path, the selected levels pass through inverse quantization
and the inverse transform before prediction is added. For INTER output,
\[
  X[y,x]=\operatorname{clip}_{[0,255]}\big(P[y,x]+D[y,x]\big).
\]
The inverse result $D$ lies in $[-256,255]$, and $P$ lies in $[0,255]$, so their
sum lies in $[-256,510]$. Signed halfwords accommodate the addition. Clip the
complete sum, then store bytes. Clipping a negative residual to zero first would
erase a needed correction. Adding in bytes would wrap before clipping.

\begin{table}[!ht]
  \centering
  \begin{tabular}{l r l}
    \toprule
    Boundary             & First value & Meaning                       \\
    \midrule
    Current / prediction & 130 / 128   & Unsigned samples              \\
    Transform input      & 2           & Signed residual               \\
    Forward DC           & 16          & Coefficient                   \\
    Selected DC level    & 1           & Ordinary INTER level          \\
    Reconstructed DC     & 23          & Inverse-quantized coefficient \\
    Inverse output       & 3           & Spatial correction            \\
    Reconstruction       & 131         & Future prediction sample      \\
    \bottomrule
  \end{tabular}
  \caption{The constant-block trace uses the worked transform and $Q=8$.}
\end{table}

The ideal DC-only inverse for coefficient 23 is $23/8=2.875$ at each sample.
The supplied transform rounds it to three. The current block was 130, but its
reconstruction is 131. Now let the next raw block be 132. Prediction from 131
leaves residual one and DC eight, which the same quantizer maps to zero. Its
reconstruction remains 131. Both encoder and decoder can agree while differing
from the raw source.

If the encoder instead predicts from raw 130, it selects residual two, level one
and inverse correction three again. Its faulty local result becomes 133, while
the decoder adds three to retained 131 and obtains 134. That disagreement is
additional to quantization loss. Repeating an unchanged raw input picture does
not by itself make skipping correct: the reference may still differ from it.

The block encoder exposes selected levels, reconstructed coefficients, signed
inverse output and final bytes in one 448-byte result. Inspect those boundaries
in order. Final clipping can hide an incorrect intermediate because several
out-of-range sums become the same endpoint. Keep the last accepted picture
separate from the new reconstruction until the picture and required output have
succeeded. The next chapter extends that publication rule across stream writes.

\begin{figure}[!ht]
\centering
\begin{tikzpicture}[alt={Four rows of native observation pairs show first level one and reconstructed DC 23; first inverse output three and sample 131; next level zero and reconstructed DC zero; next inverse output zero and sample 131. The two retained sample values have outlines.}]
\node[anchor=north west,inner sep=0,font=\small\bfseries] (title) at (0,0) {Follow selected levels into the next reference};
\node[anchor=north west,inner sep=0] (a) at ([yshift=-2mm]title.south west) {\MARSRevisionCrop{revision/captures/h263-prediction/h263-chain.png}{0.99\linewidth}{0.03125}{0.26888889}{0.58203125}{0.60888889}};
\begin{scope}[shift={(a.south west)},x={($(a.south east)-(a.south west)$)},y={($(a.north west)-(a.south west)$)}]
\draw[black,line width=.7pt] (0.582500,0.386818) rectangle (0.965000,0.569727);
\draw[black,dashed,line width=.7pt] (0.582500,0.005000) rectangle (0.965000,0.187909);
\end{scope}
\end{tikzpicture}
\caption{The caller copies four pairs from its two actual block results: selected DC and reconstructed DC, then inverse output and sample, first for raw 130 and then for raw 132. The outlined and dashed sample fields are both 131. The second call uses the first result as its prediction.}
\label{fig:revision-h263-native-chain}
\end{figure}


\Needspace{30\baselineskip}
\section{Implement and explain one reconstruction boundary}

The independent \texttt{h263-reconstruct8.asm} accepts predicted bytes and signed
inverse halfwords for an eight-by-eight block. It publishes clipped sample bytes
and returns how many sums lay strictly outside zero through 255. A sum exactly
zero or 255 is already representable and does not count as clipped.

The demo predicts 128 throughout. Its correction row is
$(-129,-128,-127,0,126,127,128,255)$, so the output row is
$(0,0,1,128,254,255,255,255)$. Three sums in each row need clipping, giving 24
for eight rows. This exercise isolates signed addition, publication and access
bounds; it does not choose coefficients or replace the inverse-accuracy gate.

Validate the strides, alignments and address extents before reading samples.
Then admit all corrections in the inverse range before publishing any result.
One solution validates then walks again. Another stages sixty-four bytes privately
and publishes them after the last sample passes. Both can support output exactly
aliasing prediction when their strides agree. A late invalid correction must leave
that aliased prediction unchanged. The lab gives the complete interface, feedback
command and questions for testing these choices independently.

\begin{figure}[!ht]
\centering
\begin{tikzpicture}[alt={The completed independent reconstruction demonstration returns status zero and clipped-sample count 24, copied into native observation fields. Count 24 is outlined.}]
\node[anchor=north west,inner sep=0,font=\small\bfseries] (title) at (0,0) {Count sums that require clipping};
\node[anchor=north west,inner sep=0] (a) at ([yshift=-2mm]title.south west) {\MARSRevisionCrop{revision/captures/h263-prediction/h263-reconstruction.png}{0.99\linewidth}{0.03125}{0.33888889}{0.58203125}{0.60888889}};
\begin{scope}[shift={(a.south west)},x={($(a.south east)-(a.south west)$)},y={($(a.north west)-(a.south west)$)}]
\draw[black,line width=.7pt] (0.582500,0.005000) rectangle (0.965000,0.443809);
\end{scope}
\end{tikzpicture}
\caption{The completed demonstration returns zero and 24 clipped sums. The first and last representable endpoints themselves are not counted. The caller copies these return values before its ordinary status/count fields are stored.}
\label{fig:revision-h263-native-reconstruction}
\end{figure}


\section*{Sources and next connection}

The selected baseline interpolation, inverse coefficient meaning, inverse transform
and sample reconstruction follow clauses 6.1--6.3 and Annex A of
\href{https://www.itu.int/rec/T-REC-H.263-199802-S/en}{ITU-T H.263 (02/1998)}.
The search radius, tie policy and forward quantizer are project choices.
\path{docs/h263-profile.md} records the restrictions and complete assignment.
The following chapter must encode decisions as legal syntax and verify the
result using an independent decoder; correct numerical kernels alone cannot
establish an interoperable video stream.
